% Chapter 4 — Dynamic Characterization
% Converted from thesis_notes/draft_ch4_dynamic.md
\chapter{Dynamic Characterization}
\label{chap:dynamic}
\lhead{\emph{Dynamic Characterization}}

\section{Objective}
\label{sec:dyn-obj}

\cref{chap:static} fixed the steady-state map; this chapter identifies how the motor moves \emph{between} operating points. The transient is modelled as a first-order-plus-dead-time (FOPDT) response~\citep{seborg2016}, and the central questions are structural: is one global time constant enough, or does the dynamic behaviour depend on operating point, on direction, and on whether the motor is speeding up or slowing down? The answers define the working model---termed \textbf{Model~C}---as a cascade in which a static block (\cref{chap:static}) feeds a dynamic block whose parameters are indexed by direction and by acceleration/deceleration region.

\section{Method}
\label{sec:dyn-method}

\textbf{Trials.} Step responses were collected with the automated sequencer firmware (\sref{sec:firmware}), which removes operator timing variability. Acceleration (``step-up'') trials drive the motor from rest to a target PWM; deceleration (``step-down'') trials hold PWM~=~240 to a settled speed and then step down to a target. The Phase~2.1 session covered PWM $\in \{160, 200, 240\}$ accel ($\times 3$ runs) and $240\to\{0, 160\}$ decel ($\times 3$ runs); the gap-fill session added PWM $\in \{180, 220\}$ accel ($\times 5$ runs) and $240\to\{100, 180, 220\}$ decel ($\times 3$ runs), plus eight piggyback re-runs of Phase~2.1 conditions (\sref{sec:xsession}). The piggyback re-runs are not folded into Tables~\ref{tab:accel}--\ref{tab:decel}; they are reported separately as the cross-session drift signal and quantified in \cref{chap:validation}.

\textbf{Model.} Acceleration is fitted as a rise from rest,
\begin{equation}
  \mathrm{rpm}(t) = \Kss\left(1 - e^{-(t - t_0 - \Td)/\tau}\right),
  \label{eq:accel}
\end{equation}
and deceleration as a decay between two levels,
\begin{equation}
  \mathrm{rpm}(t) = \yfinal + (\yinit - \yfinal)\,e^{-(t - t_0 - \Td)/\tau},
  \label{eq:decel}
\end{equation}
with gain $\Kss$ (or the level pair $\yinit, \yfinal$), dead time $\Td$, and time constant $\tau$.

\textbf{Fitting.} Each trial is fitted in two stages: a 63.2\,\%-rise heuristic provides robust initial estimates, then a nonlinear least-squares solver refines them under physical bounds~\citep{ljung1999}. The bounds enforce $\tau \in [10, 1000]$\,ms and $\Td \in [0, 200]$\,ms; the gain is sign-constrained by direction with magnitude up to twice the heuristic estimate (and, for deceleration, the level pair bounded to $\pm 400$\,rpm), while $t_0$ is fixed at the step-command index. The heuristic alone snaps $\Td$ and $\tau$ to the 10\,ms sample grid; the nonlinear fit yields continuous parameters and per-window residuals (rise, steady-state, total). All fits converged---18/18 and 12/12 in Phase~2.1, 20/20 and 18/18 in gap-fill. Residual windows are recorded so that validation (\cref{chap:validation}) can compare like with like: the rise window for accel, and the post-2000\,ms fit window for decel.

\section{Acceleration dynamics}
\label{sec:accel}

\begin{table}[t]
  \centering
  \caption{Acceleration FOPDT parameters, all ten operating points (mean $\pm$ std). $\Td$ is an effective curve-fit parameter (finding~4), not a physical latency.}
  \label{tab:accel}
  \small
  \begin{tabular}{@{}llrrcrrl@{}}
    \toprule
    Dir & PWM & $\Kss$ (rpm) & $\Td$ (ms) & $\tau$ (ms) & rms$_{\mathrm{rise}}$ & $n$ & Session \\
    \midrule
    fwd & 160 & 52.08    & ${\approx}0$ & $200.1 \pm 22.0$ & 2.61 & 3 & P2.1 \\
    fwd & 180 & 89.55    & 2.33  & $102.9 \pm 6.87$ & 3.70 & 5 & gap-fill \\
    fwd & 200 & 123.53   & 4.92  & $95.1 \pm 4.64$  & 2.48 & 3 & P2.1 \\
    fwd & 220 & 175.4    & 6.43  & $108.4 \pm 2.05$ & 4.97 & 5 & gap-fill \\
    fwd & 240 & 209.74   & 17.06 & $152.3 \pm 2.74$ & 4.69 & 3 & P2.1 \\
    rev & 160 & $-69.55$ & ${\approx}0$ & $191.7 \pm 12.15$ & 4.03 & 3 & P2.1 \\
    rev & 180 & $-93.97$ & 1.25  & $93.08 \pm 0.88$ & 3.35 & 5 & gap-fill \\
    rev & 200 & $-131.35$& 5.63  & $81.24 \pm 0.94$ & 2.66 & 3 & P2.1 \\
    rev & 220 & $-180.0$ & 7.47  & $102.6 \pm 1.57$ & 5.10 & 5 & gap-fill \\
    rev & 240 & $-222.56$& 23.49 & $144.4 \pm 2.45$ & 7.17 & 3 & P2.1 \\
    \bottomrule
  \end{tabular}
\end{table}

The acceleration FOPDT parameters are collected in \tref{tab:accel}; four findings follow.

\textbf{(1) $\tau$ is operating-point dependent---a broad bathtub.} The time constant is smallest in the mid-range and rises at both ends: $\tau \approx 81$--$110$\,ms across PWM 180--220, climbing to ${\approx}145$--$152$\,ms at PWM~=~240 and ${\approx}192$--$200$\,ms at PWM~=~160 (\fref{fig:accel_tautd}, \fref{fig:accel_tau}). The bottom of the curve is flat across a ${\approx}40$-PWM window rather than a single sharp minimum---an earlier ``V-shape'' reading is refined here to a \textbf{bathtub}. A single global $\tau$ (call it ``Model~B'') is therefore misspecified by roughly $2\times$ across the range; a region-dependent dynamic block (Model~C) is required. The rise at PWM~=~160 is friction-dominated (the operating point sits just above breakaway); the rise at PWM~=~240 reflects departure from first-order behaviour near full drive (back-EMF/saturation), evidenced by the residuals below.

\begin{figure}[t]
  \centering
  \includegraphics[width=0.95\textwidth]{07_tau_td_vs_pwm}
  \caption{Acceleration time constant $\tau$ and (effective) dead time $\Td$ versus PWM, forward and reverse (Phase~2.1, mean $\pm$ std). $\tau$ rises at both ends of the operating range.}
  \label{fig:accel_tautd}
\end{figure}

\begin{figure}[t]
  \centering
  \includegraphics[width=0.95\textwidth]{16_phase2gapfill_accel_tau_vs_pwm}
  \caption{Acceleration $\tau$ and $|\Kss|$ versus PWM, combining Phase~2.1 (160/200/240) and gap-fill (180/220): the broad bathtub minimum spans PWM 180--220, while $|\Kss|$ stays linear in PWM.}
  \label{fig:accel_tau}
\end{figure}

\textbf{(2) The gain is linear and direction-clean.} $\Kss$ rises smoothly and almost linearly with PWM within each direction, with very tight run-to-run spread (per-condition std $\le 2$\,rpm across the five gap-fill runs). The static-block separability of \cref{chap:static} thus holds across all five calibrated PWMs---the dynamics change with operating point, but the steady gain does not misbehave.

\textbf{(3) The dynamic block is direction-symmetric to first order.} At every PWM, $\tau_{\mathrm{fwd}}$ and $\tau_{\mathrm{rev}}$ agree to within a small margin, with reverse running consistently a few percent faster. The largest gap is at PWM~=~200 (fwd 95\,ms vs rev 81\,ms, ${\approx}15\,\%$), which exceeds the run-to-run scatter and is flagged as a genuine but minor effect---possible static/dynamic coupling or commutation asymmetry---that does not undermine the cascade. Crucially, the \emph{large} directional asymmetry of \cref{chap:static} does \textbf{not} reappear in the dynamics, confirming that asymmetry lives in the static block.

\textbf{(4) $\Td$ is an effective parameter, not physical latency.} The fitted dead time is ${\approx}0$ at PWM~=~160 and inflates to 17--23\,ms at PWM~=~240---the opposite of a fixed transport delay. This is the FOPDT model absorbing rise-shape error at the operating-point extremes; the genuine encoder/serial latency (${\approx}10$\,ms, one sample) is what the heuristic recovers. $\Td$ in \tref{tab:accel} should therefore be read as a curve-fit nuisance parameter, and PWM~=~200---where the rise is cleanly first-order (rms ${\approx}2.5$\,rpm)---taken as the canonical operating point. Fit quality degrades away from it: rms$_{\mathrm{rise}}$ climbs to ${\approx}4.7$\,rpm (PWM 240 fwd) and ${\approx}7.2$\,rpm (PWM 240 rev), with the residual grid (\fref{fig:accel_resid}) showing structured, not random, error there.

\begin{figure}[t]
  \centering
  \includegraphics[width=0.95\textwidth]{08_fopdt_curvefit_residuals}
  \caption{Curve-fit FOPDT rises with residual panels for the six Phase~2.1 acceleration conditions. Residuals are random at PWM~=~200 but structured at PWM~=~160 (friction onset) and PWM~=~240 (smooth onset), confirming effective-FOPDT behaviour at the extremes.}
  \label{fig:accel_resid}
\end{figure}

\section{Deceleration dynamics}
\label{sec:decel}

\begin{table}[t]
  \centering
  \caption{Deceleration FOPDT parameters, all ten transitions from PWM~=~240 (mean $\pm$ std).}
  \label{tab:decel}
  \footnotesize
  \begin{tabular}{@{}lrrrrcrrl@{}}
    \toprule
    Dir & $240\to$ & $\yinit$ & $\yfinal$ & $\Td$ (ms) & $\tau$ (ms) & rms$_{\mathrm{dec}}$ & $n$ & Session \\
    \midrule
    fwd & 0   & 214.3   & $-0.69$  & 55.19 & $172.2 \pm 1.77$  & 8.67 & 3 & P2.1 \\
    fwd & 100 & 220.5   & 1.09     & 25.51 & $366.2 \pm 1.62$  & 4.74 & 3 & gap-fill \\
    fwd & 160 & 215.7   & 73.75    & 14.06 & $275.3 \pm 8.77$  & 4.39 & 3 & P2.1 \\
    fwd & 180 & 219.9   & 94.18    & 0.33  & $356.6 \pm 3.55$  & 4.28 & 3 & gap-fill \\
    fwd & 220 & 219.4   & 175.9    & 8.92  & $122.9 \pm 16.14$ & 5.43 & 3 & gap-fill \\
    rev & 0   & $-229.1$& ${\approx}0$ & 62.93 & $189.5 \pm 3.87$ & 9.78 & 3 & P2.1 \\
    rev & 100 & $-229.9$& $-0.61$  & 25.27 & $349.0 \pm 20.43$ & 5.58 & 3 & gap-fill \\
    rev & 160 & $-231.7$& $-85.66$ & 9.95  & $306.4 \pm 3.99$  & 2.82 & 3 & P2.1 \\
    rev & 180 & $-229.8$& $-96.23$ & 5.94  & $321.3 \pm 19.13$ & 4.33 & 3 & gap-fill \\
    rev & 220 & $-230.2$& $-181.8$ & 11.43 & $109.2 \pm 5.13$  & 5.49 & 3 & gap-fill \\
    \bottomrule
  \end{tabular}
\end{table}

The deceleration parameters (\tref{tab:decel}) yield four further findings.

\textbf{(1) Deceleration is slower than acceleration.} In the same operating range, decel $\tau$ exceeds every accel $\tau$: $240\to160$ decays with $\tau \approx 275$\,ms (fwd) / 306\,ms (rev), against accel $\tau$ of 95--152\,ms at PWM 200--240. Speeding up and slowing down are not the same process, which is the primary justification for splitting Model~C into separate accel and decel regions.

\textbf{(2) Decel $\tau$ is non-monotone in the target.} It is \emph{not} simply ``slower for bigger steps.'' The slowest decays are to targets sitting on the deadzone edge---$240\to100$ and $240\to180$ ($\tau \approx 321$--$366$\,ms)---where the shaft must coast through the high-friction low-speed region. A full coast to rest ($240\to0$) is faster (172--190\,ms), and a mild decel ($240\to220$) is fastest (109--123\,ms). The deadzone, characterized statically in \cref{chap:static}, thus reshapes the \emph{dynamics} of any transition that passes near it (\fref{fig:decel_tau}, \fref{fig:decel_overlay}).

\begin{figure}[t]
  \centering
  \includegraphics[width=0.95\textwidth]{17_phase2gapfill_decel_tau_vs_target}
  \caption{Deceleration $\tau$ and $|\yfinal|$ versus target PWM (from PWM~=~240), Phase~2.1 (targets 0, 160) plus gap-fill (100, 180, 220). $\tau$ peaks for coast-through-deadzone targets.}
  \label{fig:decel_tau}
\end{figure}

\begin{figure}[t]
  \centering
  \includegraphics[width=0.95\textwidth]{18_phase2gapfill_decel_overlays}
  \caption{Deceleration trajectory overlays, Phase~2.1 versus gap-fill, at nearby targets: the persistent low-speed slowdown at deadzone-edge targets is visible directly in the traces.}
  \label{fig:decel_overlay}
\end{figure}

\textbf{(3) Regenerative braking is suppressed at non-zero targets.} Decel to 160 is slower than decel to 0 even though it is a smaller speed change. The likely cause is that at a non-zero PWM hold the L298N keeps chopping, preventing the regenerative braking that a pure coast (PWM~=~0) allows---an artefact of the driver, again localising a behavioural quirk to the static/driver side.

\textbf{(4) The direction asymmetry reverses sign between regions.} In acceleration reverse is slightly \emph{faster} (\sref{sec:accel}); in deceleration reverse is slightly \emph{slower} ($240\to160$: 306 vs 275\,ms; $240\to0$: 190 vs 172\,ms). Equal magnitude, opposite sign. This indicates partial coupling between the static and dynamic blocks---the cascade holds but is not perfectly clean, and this is carried as an explicit thesis caveat. Finally, the coast-to-rest cases ($240\to0$, $\yfinal \approx 0$) are Coulomb-friction-dominated and so only effectively first-order; their residuals (rms ${\approx}8.7$--$9.8$\,rpm over the decel window) are the largest in the set but remain bounded.

\section{Region structure: the definition of Model C}
\label{sec:modelc}

The findings compose into the working model. The plant is a cascade PWM\,$\to$\,static driver-voltage block (\cref{chap:static})\,$\to$\,motor dynamic block\,$\to$\,RPM, in which the dynamic block is \textbf{not} a single LTI system but a family of FOPDT responses indexed by:
\begin{itemize}
  \item \textbf{direction} (forward / reverse)---required by the static block, optional for the dynamics (small effect);
  \item \textbf{region} (acceleration / deceleration)---required, since decel $\tau \approx 2\times$ accel $\tau$; and
  \item \textbf{operating point}---accel $\tau$ follows a bathtub in PWM; decel $\tau$ depends non-monotonically on the target through the deadzone.
\end{itemize}

This is Model~C, the third rung of a candidate ladder fixed at the outset: \textbf{Model~A}, a single global LTI model with no static block ($u_{\mathrm{cmd}}\to$\,rpm), the straw-man baseline---quantified in \sref{sec:modelsel} and reused in \cref{chap:closedloop}; \textbf{Model~B}, the static block followed by one global dynamic block (a single gain, $\tau$, and delay); \textbf{Model~C}, the static block with a region-dependent dynamic block (this chapter); and \textbf{Model~D}, an optional friction-explicit refinement. Model~B is rejected by the ${\approx}2\times$ spread in $\tau$ across operating points; Model~C is adopted as the working model; the Coulomb/Stribeck Model~D~\citep{armstrong1994}, motivated by the PWM~=~160 and coast-to-rest residuals, is left to future work. The next section makes these dismissals quantitative.

\section{Model selection: the A/B/C ladder, quantified}
\label{sec:modelsel}

\sref{sec:modelc} adopted Model~C on structural grounds. This section makes the choice quantitative: rather than dismiss the simpler models by argument, all three are \emph{fitted to the same data} and scored head-to-head, so that the superiority of the cascade---and the separate contributions of its two ideas, the static block and the region-dependent dynamics---can be read off directly, forward and reverse. The comparison uses the Phase~2.1 grid (30 trials, single session, so that cross-session drift does not confound the structural question), and all three models are simulated with the identical FOPDT core used for validation (\cref{chap:validation}); differences are therefore due to model \emph{structure}, not numerics.

\textbf{The three models, fit on a common footing.}
\begin{itemize}
  \item \textbf{Model~A (naive global LTI).} No static block: a single proportional gain maps command to speed, $\mathrm{rpm} = K_A\,\mathrm{pwm}$, followed by one global FOPDT. Fitted over all thirty trials, $K_A = 0.785$\,rpm/PWM, $\tau = 207$\,ms, $\Td = 20$\,ms.
  \item \textbf{Model~B (static block + one global dynamic block).} The static block supplies the correct per-operating-point steady state (the deadzone and the direction-specific gain of \cref{chap:static}), but a \emph{single} global time constant governs every transition: $\tau = 187$\,ms, $\Td = 16$\,ms.
  \item \textbf{Model~C (the cascade).} The static block plus the region- and operating-point-dependent FOPDT of Tables~\ref{tab:accel}--\ref{tab:decel}, unchanged.
\end{itemize}
The pairing is deliberate: \textbf{A\,$\to$\,B isolates the value of the static block}, and \textbf{B\,$\to$\,C isolates the value of region-dependent dynamics}.

\begin{table}[t]
  \centering
  \caption{Model-ladder scoreboard on the Phase~2.1 grid: post-step RMSE (rpm) and FIT\% by region and direction, with the overall mean over all 30 trials. Model~C is best in every group, forward and reverse. Lower RMSE and higher FIT\% are better; $\mathrm{FIT\%}=0$ is no better than predicting the mean.}
  \label{tab:ladder}
  \small
  \begin{tabular}{@{}llrrrrrr@{}}
    \toprule
    & & \multicolumn{3}{c}{RMSE (rpm)} & \multicolumn{3}{c}{FIT\%} \\
    \cmidrule(lr){3-5}\cmidrule(lr){6-8}
    Region & Dir & A & B & C & A & B & C \\
    \midrule
    accel & fwd & 42.1 & 6.0 & \textbf{3.8} & $-332$ & 56.6 & \textbf{71.8} \\
    accel & rev & 38.4 & 7.2 & \textbf{4.3} & $-191$ & 53.9 & \textbf{73.4} \\
    decel & fwd & 28.3 & 6.7 & \textbf{4.3} & $-13$  & 76.4 & \textbf{85.2} \\
    decel & rev & 24.4 & 8.3 & \textbf{3.8} & $+11$  & 73.3 & \textbf{88.7} \\
    \midrule
    \multicolumn{2}{@{}l}{overall (30)} & 34.7 & 6.96 & \textbf{4.07} & $-157$ & 63.1 & \textbf{78.3} \\
    \bottomrule
  \end{tabular}
\end{table}

Each trial is scored on its post-step window by RMSE and by the normalised fit $\mathrm{FIT\%} = 100\,(1 - \lVert y - \hat{y}\rVert / \lVert y - \bar{y}\rVert)$. The scoreboard is \tref{tab:ladder}; three readings follow, and they hold in both directions.

\textbf{(1) Model~A fails outright.} Its FIT\% is \emph{negative} in three of the four groups---the naive model predicts the step responses worse than a horizontal line through their own mean. With no static block, its single gain cannot represent the deadzone, so it overshoots heavily at low PWM (at PWM~=~160 it predicts ${\approx}125$\,rpm against a measured ${\approx}52$). \fref{fig:ladder_overlay} shows the symptom directly: Model~A reaches the wrong level and rises with the wrong shape. This is the empirical justification---not the assumption---that a low-cost H-bridge drive cannot be treated as a single linear system.

\begin{figure}[t]
  \centering
  \includegraphics[width=\textwidth]{20_model_ladder_overlay}
  \caption{Model ladder on real step responses (run~01): Model~A (naive LTI), Model~B (static\,+\,global $\tau$) and Model~C (static\,+\,region $\tau$) overlaid on the measured trace for forward/reverse acceleration and deceleration. Model~A reaches the wrong level and rises with the wrong shape; Model~C tracks both level and transient. Per-panel RMSE is annotated.}
  \label{fig:ladder_overlay}
\end{figure}

\textbf{(2) The static block is the dominant improvement (A\,$\to$\,B).} Restoring the deadzone and the direction-specific steady-state map collapses overall RMSE from ${\approx}35$ to ${\approx}7$\,rpm and lifts FIT\% from $-157$ to $+63\,\%$. The single largest gain in the whole study comes from \emph{decomposing} the plant---the central claim of the cascade, here measured rather than asserted.

\textbf{(3) Region-dependent dynamics are a real, smaller gain (B\,$\to$\,C).} Model~C lowers overall RMSE further to ${\approx}4$\,rpm and FIT\% to ${\approx}78\,\%$. The improvement is modest on the full window---which is dominated by the (shared) steady state---but it concentrates exactly where \sref{sec:accel} predicts: in the \emph{transient}. Scored on a transient window (five time constants, the same window for all three models), Model~C's error is ${\approx}5.8$\,rpm against Model~B's ${\approx}15.3$\,rpm. A single global $\tau$ is misspecified by ${\approx}2\times$ across the operating range (the bathtub), and that misspecification surfaces precisely in the rise and decay, which region scheduling repairs (\fref{fig:ladder_score}).

\begin{figure}[t]
  \centering
  \includegraphics[width=\textwidth]{21_model_ladder_scoreboard}
  \caption{A/B/C scoreboard across direction\,$\times$\,region: post-step RMSE (left, lower is better) and FIT\% (right, higher is better). Model~C wins every group; FIT\% is clipped at $-50$ for the failing Model~A.}
  \label{fig:ladder_score}
\end{figure}

On the Phase~2.1 grid the candidate ladder therefore resolves monotonically and identically across direction and regime: $\mathrm{A} \ll \mathrm{B} < \mathrm{C}$. Model~C is the best-fit model forward and reverse, accelerating and decelerating, with the static block accounting for the bulk of the improvement and the region-dependent dynamics for the remainder, concentrated in the transient. This is in-sample \emph{structural} selection; the model's ability to predict held-out trials, and how far it transfers across sessions, are the separate subjects of \cref{chap:validation}.

\textbf{Extension to the full operating envelope.} Re-running the comparison across all sixty-eight calibrated trials---adding the gap-fill acceleration points (PWM 180/220) and deceleration targets ($240\to100/180/220$), with Models~A and~B re-fit globally over the enlarged set and Model~C using each condition's session-appropriate parameters---confirms the conclusion and sharpens it (\tref{tab:ladder_full}). The ordering $\mathrm{A}\ll\mathrm{B}<\mathrm{C}$ replicates \emph{independently in each session}, so it is not an artefact of a single bench day. More tellingly, the added mid-range points are exactly where Model~B is worst: its \emph{transient} error peaks at PWM 180--220 (${\approx}18$--$31$\,rpm)---the bottom of the $\tau$ bathtub, where the true time constant is fastest (${\approx}95$--$108$\,ms) and a single global $\tau$ (${\approx}211$\,ms) is most misspecified---while Model~C holds ${\approx}3$--$5$\,rpm throughout (\fref{fig:ladder_full}). Broadening the grid thus \emph{strengthens} the case for region-dependent dynamics rather than merely extending it. The consolidated grid mixes two sessions (as do Tables~\ref{tab:accel}--\ref{tab:decel}); the small extra cross-session penalty falls on all three models alike, and the single-session result of \tref{tab:ladder} remains the clean reference.

\begin{table}[t]
  \centering
  \caption{Full-envelope ladder (68 trials, Phase~2.1 $+$ gap-fill): mean post-step RMSE (rpm) and FIT\%, overall and by session. The $\mathrm{A}\ll\mathrm{B}<\mathrm{C}$ ordering replicates in both sessions.}
  \label{tab:ladder_full}
  \small
  \begin{tabular}{@{}lrrrrrr@{}}
    \toprule
    & \multicolumn{3}{c}{RMSE (rpm)} & \multicolumn{3}{c}{FIT\%} \\
    \cmidrule(lr){2-4}\cmidrule(lr){5-7}
    Subset & A & B & C & A & B & C \\
    \midrule
    Phase~2.1 ($n{=}30$) & 34.8 & 7.1 & \textbf{4.1} & $-157$ & 62.8 & \textbf{78.3} \\
    gap-fill ($n{=}38$)  & 35.2 & 9.7 & \textbf{4.5} & $-115$ & 43.1 & \textbf{67.6} \\
    \midrule
    overall ($n{=}68$)   & 35.0 & 8.5 & \textbf{4.3} & $-133$ & 51.8 & \textbf{72.3} \\
    \bottomrule
  \end{tabular}
\end{table}

\begin{figure}[t]
  \centering
  \includegraphics[width=\textwidth]{22_model_ladder_fullgrid}
  \caption{Full operating envelope: transient RMSE versus operating point for the three models (log scale; forward and reverse averaged). Model~B's single global $\tau$ fails worst at the acceleration bathtub bottom (PWM 180--220, shaded) and at deadzone-edge deceleration targets, whereas Model~C stays flat-low across the envelope.}
  \label{fig:ladder_full}
\end{figure}

\section{Cross-session note}
\label{sec:xsession}

Tables~\ref{tab:accel} and~\ref{tab:decel} combine two bench sessions two weeks apart: PWM 160/200/240 and targets 0/160 from Phase~2.1, and PWM 180/220 and targets 100/180/220 from the gap-fill session (the \textsf{Session} column). The bathtub and target-dependence framings rely on combining both. Within either session, reproducibility is excellent (run-to-run $\tau$ std mostly $\le$ a few ms). Between sessions there is structured drift---most visibly, the gap-fill reverse high-PWM $\tau$ runs ${\approx}20\,\%$ below its Phase~2.1 value (${\approx}114$ vs 144\,ms in the piggyback re-runs), and low-speed-target decel runs slower than Phase~2.1 interpolation would predict. The consolidated tables are therefore an honest two-session composite, and the magnitude and structure of the drift are quantified as a validation result in \cref{chap:validation} rather than smoothed over here.

\section{Summary}
\label{sec:dyn-summary}

The motor's transient response is well captured, operating point by operating point, as FOPDT. The dynamics are operating-point dependent (accel $\tau$ a broad bathtub, 81--200\,ms), region dependent (decel $\tau \approx 2\times$ accel, and non-monotone through the deadzone), and only weakly direction dependent---with the direction asymmetry switching sign between accel and decel, a small but real departure from a perfectly clean cascade. The fitted dead time is an effective parameter; PWM~=~200 is the canonical clean FOPDT point, and the extremes are effective fits with bounded, structured residuals that motivate a future friction-explicit model. The twenty-condition parameter set of Tables~\ref{tab:accel} and~\ref{tab:decel} is the dynamic block of Model~C; \cref{chap:validation} tests how well it predicts held-out trials and how far it transfers across runs and sessions.

\chaptersources{\texttt{data/processed/phase21\_fopdt\_curvefit\_per\_condition.csv}, \texttt{phase2gapfill\_fopdt\_accel\_per\_condition.csv} (\tref{tab:accel}) \textperiodcentered{} \texttt{phase21\_stepdown\_curvefit\_per\_condition.csv}, \texttt{phase2gapfill\_stepdown\_per\_condition.csv} (\tref{tab:decel}) \textperiodcentered{} \texttt{notebooks/03\_step\_responses.ipynb}, \texttt{05\_gap\_fill.ipynb} (fits, figures 05--11, 16--18) \textperiodcentered{} status-doc \S4.2--\S4.3, \S5 \textperiodcentered{} \texttt{thesis\_notes/log.md} 2026-05-07 / 2026-05-22 (interpretation, accel/decel asymmetry, bathtub refinement).}
